\documentclass[sigconf,9pt,nonacm,balance=false]{acmart}

\usepackage{enumitem}
\usepackage{listings}
\usepackage{xurl}
\IfFileExists{libertine.sty}{}{\usepackage[T1]{fontenc}\usepackage{lmodern}}
\urlstyle{tt}
\settopmatter{printacmref=false,authorsperrow=3}
\setlist[itemize]{leftmargin=*,nosep}
\lstset{basicstyle=\ttfamily\footnotesize,columns=fullflexible,
  keepspaces=true,breaklines=true,showstringspaces=false,
  aboveskip=5pt,belowskip=5pt}
\newcommand{\repo}{https://github.com/ddps-lab/KubePACS}
\newcommand{\pathcode}[1]{\nolinkurl{#1}}

\begin{document}

\title[Artifact Appendix for KubePACS]{Artifact Appendix for KubePACS:
Kubernetes Cluster Using Performant, Highly Available,
and Cost Efficient Spot Instances}

\author{Taeyoon Kim}
\authornote{Both authors contributed equally to this work.}
\email{tykim7@hanyang.ac.kr}
\orcid{0009-0003-0705-0713}
\affiliation{\department{Dept. of Data Science}
  \institution{Hanyang University}\city{Seoul}\country{Republic of Korea}}
\author{Kyumin Kim}
\authornotemark[1]
\email{okkimok123@hanyang.ac.kr}
\orcid{0009-0007-2705-1122}
\affiliation{\department{Dept. of Data Science}
  \institution{Hanyang University}\city{Seoul}\country{Republic of Korea}}
\author{Enrique Molina-Gim\'enez}
\email{enrique.molina@urv.cat}
\orcid{0009-0005-2597-3815}
\affiliation{\department{Dept. of Computer Eng. and Math.}
  \institution{Universitat Rovira i Virgili}\city{Tarragona}\country{Spain}}
\author{Pedro Garc\'ia-L\'opez}
\email{pedro.garcia@urv.cat}
\orcid{0000-0002-9848-1492}
\affiliation{\department{Dept. of Computer Eng. and Math.}
  \institution{Universitat Rovira i Virgili}\city{Tarragona}\country{Spain}}
\author{Kyungyong Lee}
\authornote{Corresponding author.}
\email{kyungyong@hanyang.ac.kr}
\orcid{0000-0003-0312-4386}
\affiliation{\department{Dept. of Data Science}
  \institution{Hanyang University}\city{Seoul}\country{Republic of Korea}}

\renewcommand{\shortauthors}{Kim et al.}
\begin{abstract}
This artifact accompanies KubePACS, a Kubernetes-native spot-instance
provisioning system that jointly considers hardware performance, spot
pricing, and multi-node availability. It provides the optimizer,
Kubernetes integration sources, stored experimental results,
and scripts for generating the paper's figures and tables. To minimize cloud
cost and setup effort, evaluation runs the optimizer on stored inputs and
generates figures and tables from the supplied results locally. We request the
\emph{Artifacts Available}, \emph{Artifacts Functional}, and
\emph{Results Reproduced} badges.
\end{abstract}
\maketitle

\section{Artifact Check-List}
\begin{itemize}
  \item \textbf{Programs:} Python analysis and result-generation scripts with
    accompanying optimizer and Kubernetes integration sources.
  \item \textbf{Software:} Python 3.11 and \texttt{uv} for figures, using the
    supplied lock file. The optimizer check adds PuLP~2.9.0 and
    requests~2.32.5. Table generation uses the Python standard library.
  \item \textbf{Run-time environment:} Linux workstation for local analysis.
    Internet access for repository retrieval and dependency installation.
  \item \textbf{Hardware:} General-purpose CPU. No GPU required for local
    result generation.
  \item \textbf{Disk space:} Approximately 4 GiB for the local figure workflow,
    including dependencies, temporary inputs, and generated outputs.
  \item \textbf{Execution time:} Allow several minutes for all figures.
    Each plotting script has a 180-second timeout. Allow several additional
    minutes for the optimizer check.
  \item \textbf{Outputs:} 24 figure PDFs and Tables 2 and 3 in CSV, Markdown,
    and LaTeX. Optimizer recommendations are saved as three JSON reports.
    Figure execution logs and a machine-readable summary are also provided.
  \item \textbf{Access:} Source, data, and instructions at \url{\repo}.
\end{itemize}

\section{Description}
\subsection{Access and Structure}
The repository is the entry point for installation and evaluation. Commands
in this appendix start at its root.
\begin{lstlisting}
KubePACS/
|-- README.md
|-- figures/
|   |-- reproduce.py
|   |-- figure*/
|   |-- table2_alpha/
|   +-- table3_compute/
+-- api/
\end{lstlisting}

\subsection{Software and Data}
\textbf{Optimizer and integration.}
The optimizer combines integer linear programming with Golden Section
Search to select instance types, availability zones, and node counts.
The \pathcode{api/} directory exposes this implementation through a Lambda
handler that can also be invoked locally. Kubernetes integration sources
are also supplied.

\textbf{Experimental inputs.}
The \pathcode{figures/} directory contains benchmark data, stored simulation
and application results, plotting scripts, and reference PDFs. Historical
price inputs are supplied locally, so figure generation does not query live
market data. Table~2 aggregates allocation and alpha-sweep results, including
the supplied alpha=1 records. Table~3 uses the recorded compilation and video
encoding throughput and instance prices. Figures~3 and~4 are architectural
illustrations supplied as PPTX sources. Table~1 defines mathematical notation.

\section{Relationship to the Paper}
The artifact connects the paper's claims to implementations and stored
results as follows. Figure and table numbers below refer to the main paper.
\begin{itemize}
  \item \textbf{Cost--performance optimization:} benchmark and price
    comparisons and baseline evaluations in Figures~1, 5, and~10.
  \item \textbf{Multi-node availability:} recorded request fulfillment
    and availability indicators in Figures~2 and~9.
  \item \textbf{Adaptive optimization:} alpha sensitivity, search overhead,
    and configuration comparisons in Figures~6--7 and Table~2.
  \item \textbf{Workload-aware selection:} specialized instance selection
    in Figure~8.
  \item \textbf{Application efficiency:} recorded infrastructure metrics,
    graph-analytics execution times, and compute-workload comparisons in
    Figures~10--11 and Table~3.
\end{itemize}
The local scripts generate Figures~1, 2, and~5--12 and Tables~2--3
from the supplied experimental results.

\section{Installation and Execution}
\subsection{Local Figure and Table Generation}
Install Git and curl on Linux. Install \texttt{uv} and Python~3.11,
then obtain the repository and run:
\begin{lstlisting}
curl -LsSf https://astral.sh/uv/install.sh | sh
export PATH="$HOME/.local/bin:$PATH"
uv python install 3.11
git clone https://github.com/ddps-lab/KubePACS.git
cd KubePACS
uv sync --locked --project figures --python 3.11
uv run --locked --project figures \
  python figures/reproduce.py
uv run --locked --project figures \
  python figures/table2_alpha/table_02.py
uv run --locked --project figures \
  python figures/table3_compute/table_03.py
\end{lstlisting}
Dependency installation requires Internet access. The generation commands
then operate on local inputs.
Choose a new \texttt{-{}-output} directory for subsequent runs.

\textbf{Data processing and collection.}
The figure runner executes 16 scripts in a temporary input copy, leaving
reference files unchanged. It collects PDFs, per-script logs, and
\pathcode{summary.json} under \pathcode{artifact-results/figures/}.
Table outputs are collected under \pathcode{artifact-results/table2/} and
\pathcode{artifact-results/table3/}. The table commands emit
\pathcode{table2.csv/.md/.tex} and \pathcode{table3.csv/.md/.tex}. Table~2 also
emits \pathcode{table2_details.csv} with 1,200 normalized observations.

\textbf{Expected outputs.}
The figure command exits with code~0 and reports
\texttt{PASS: 16 scripts, 24 PDFs}. The summary lists script exit codes and
output-presence checks. The table commands report \texttt{PASS: Table 2}
and \texttt{PASS: Table 3}. Table~2 checks all five rounded values against the
paper: 0.8616, 0.9563, 0.0006, 0.0001, and 1.0000 for Greedy, alpha=0,
alpha=0.5, alpha=1, and GSS, respectively.

\textbf{Interpretation.}
Compare generated plots with the supplied references using the metric
definitions and figure-by-figure mapping in \pathcode{figures/README.md}.
Table~2 normalizes each configuration's efficiency to GSS for the same run
and scenario, then averages 240 ratios per configuration. Table~3 summarizes
the two compute workloads and their cost and throughput comparisons.
The local workflow reprocesses stored experimental results for generation
and inspection of the figures and tables.

\subsection{Local Optimizer Check}
Run \pathcode{figures/common/library/check_optimizer.py} using the command
in its accompanying README to invoke \texttt{getGoldenNodepool} with the
supplied merged CSV and cached AZ mapping. Both input files are included
in the repository.
These files provide prices, availability limits, hardware specifications,
benchmark scores, and zone names. The command uses the existing optimizer
with three requirements: 10 pods at 1 vCPU and 2 GiB per pod, 50 pods at
2 vCPU and 4 GiB, and 100 pods at 1 vCPU and 8 GiB.

Use repeated \texttt{-{}-case PODS,VCPU,GIB} arguments to supply different
requirements and \texttt{-{}-output} to select a fresh output directory.
Success is \texttt{PASS: 3 local optimizer scenarios}, three scenario JSON
reports, and \pathcode{summary.json} under \pathcode{artifact-results/optimizer/}.
Each report contains instance
types, zone names, instance counts, estimated hourly cost, performance score,
selected alpha, and total pod capacity. The command checks that capacity
covers the requested pod count and instance counts respect the stored
availability limits. Compare recommendations across the three requirements.

\end{document}
